\section{Conclusion And Future Work}
\label{sec:conclusion}

In this paper, we have presented a novel unsupervised framework for invariance induction in neural networks. Our method models invariance as an information separation task achieved by competitive training between a predictor and a decoder coupled with disentanglement. We described an adversarial instantiation of this framework and provided analysis of its working. Experimental evaluation shows that our unsupervised adversarial invariance induction model outperforms state-of-the-art methods, which are supervised, on learning invariance to inherent nuisance factors, effectively using synthetic data augmentation for learning invariance, and domain adaptation. Furthermore, the fact that our framework requires no annotations for variations of nuisance factors, or even knowledge of such factors, shows the conceptual superiority of our approach compared to previous methods. Since our model does not make any assumptions about the data, it can be applied to any supervised learning task, eg., binary/multi-class classification or regression, without loss of generality.

The proposed approach is not designed to learn ``fair representations'' of data, e.g., making predictions about the savings of a person invariant to age, when such bias exists in data and making the prediction task invariant to such biasing factors is of higher priority than the prediction performance~\cite{bib:sai}. In future work, we will augment our model with the capability to \emph{additionally} use supervised information (when available) about known nuisance factors for learning invariance to them, which will, consequently, help our model learn fair representations.